\documentclass{article}
\usepackage[UTF8]{ctex}
\usepackage{amsmath,amssymb,amsfonts}
\usepackage{enumitem}
\usepackage{geometry}
\geometry{a4paper, margin=1in}
% === 错题本排版优化（不参与题目标识） ===
\setlength{\parindent}{0pt}
\setlength{\parskip}{0.35em plus 0.10em minus 0.08em}
\setlength{\emergencystretch}{3em}
\setlist{itemsep=0.25em, topsep=0.40em, parsep=0pt}
\allowdisplaybreaks[2]
% === 错题本排版优化结束 ===
\begin{document}
\section*{河海大学2025年高等代数真题}
\section*{87．河海大学2025年高等代数试卷}
一．填空题．每空4分，共20分．\\

\subsection{题目1}
1．\(a, b, c, d\) 是给定的实数，在 \(\mathbb{R}^{2 \times 2}\) 上定义线性变换

\[
\mathcal{A}(X)=\left(\begin{array}{ll}
a & b \\
c & d
\end{array}\right) X, \forall X \in \mathbb{R}^{2 \times 2} .
\]

\(\mathbb{R}^{2 \times 2}\) 是 2 阶实方阵全体，它的一组基为

\[
E_{11}=\left(\begin{array}{ll}
1 & 0 \\
0 & 0
\end{array}\right), E_{12}=\left(\begin{array}{ll}
0 & 1 \\
0 & 0
\end{array}\right), E_{21}=\left(\begin{array}{ll}
0 & 0 \\
1 & 0
\end{array}\right), E_{22}=\left(\begin{array}{ll}
0 & 0 \\
0 & 1
\end{array}\right) .
\]

\(\mathcal{A}\) 在这组基下的矩阵为 \(\_\_\_\_\)．

\subsection{题目2}
2．\(V\) 是有限维线性空间， \(\mathcal{A}, \mathcal{B}\) 是 \(V\) 上的两个线性变换， \(\mathcal{A}^{2}=\mathcal{A}, \mathcal{B}^{2}=\mathcal{B}\) ，且 \((\mathcal{A}+\mathcal{B})^{2}=\mathcal{A}+\mathcal{B}\) ，那么 \(\mathcal{A} \mathcal{B}\) \(\_\_\_\_\)\(\mathcal{O}, \mathcal{O}\) 是 \(V\) 上的零变换（填＂等于＂或＂不一定等于＂）。

\subsection{题目3}
3．\(\left(\begin{array}{lll}1 & 0 & 0 \\ 0 & 1 & 1 \\ 0 & 0 & 1\end{array}\right)^{100}\left(\begin{array}{lll}1 & 2 & 3 \\ 4 & 5 & 6 \\ 7 & 8 & 9\end{array}\right)\left(\begin{array}{lll}0 & 0 & 1 \\ 0 & 1 & 0 \\ 1 & 0 & 0\end{array}\right)^{200}=\)

\subsection{题目4}
4．写出下列所有正确命题的序号 \(\_\_\_\_\)。\\
（a）\(p(x)\) 是有理数域上次数大于 0 的多项式，若对有理数域上任意两个多项式 \(f(x), g(x)\) ，由 \(p(x) \mid f(x) g(x)\) 能得出 \(p(x) \mid g(x)\) 或 \(p(x) \mid f(x)\) ，则 \(p(x)\) 是不可约多项式。\\
（b）\(f(x)\) 是有理数域上次数大于 0 的首一多项式，\(f(x)\) 是不可约多项式的方幂的充要条件为：对有理数域上任意一个多项式 \(g(x)\) ，有 \((f(x), g(x))=1\) 或者存在正整数 \(m\) 使得 \(f(x) \mid(g(x))^{m}\) ．\\
（c）\(f(x)\) 是有理数域上次数大于 0 的首一多项式，\(f(x)\) 是不可约多项式方幂的充要条件为：对有理数域上任意两个多项式 \(g(x), h(x)\) ，有 \(f(x) \mid g(x)\) 或者存在正整数 \(m\) ，使得 \(f(x) \mid(h(x))^{m}\) ．

\subsection{题目5}
5．\(a, b\) 为实数，\(A=\left(\begin{array}{cccc}3 & -2 & a & -16 \\ 2 & -3 & 0 & 1 \\ 1 & -1 & 1 & -3 \\ 3 & b & 1 & -2\end{array}\right)\) ，问 \(a, b\) 满足 \(\_\_\_\_\)时，\(A\) 的秩最小。

二、计算题．\\

\subsection{题目6}
6．（10 分）求行列式

\[
D_{n}=\left|\begin{array}{cccc}
x_{1}-m & x_{2} & \cdots & x_{n} \\
x_{1} & x_{2}-m & \cdots & x_{n} \\
\vdots & \vdots & & \vdots \\
x_{1} & x_{2} & \cdots & x_{n}-m
\end{array}\right|
\]

\subsection{题目7}
7．（10 分）\(a, b\) 为实数，问 \(a, b\) 满足什么条件时线性方程组

\[
\left\{\begin{array}{l}
a x_{1}+x_{2}+x_{3}=4 \\
x_{1}+b x_{2}+x_{3}=3 \\
x_{1}+2 b x_{2}+x_{3}=4
\end{array}\right.
\]

有唯一解，有无穷解，无解？有解时求方程组的解。

\subsection{题目8}
8．（10 分）\(a_{1}, \cdots, a_{n}\) 为 \(n\) 个实数，求

\[
A=\left(\begin{array}{cccc}
0 & & & -a_{n} \\
1 & 0 & & -a_{n-1} \\
& 1 & 0 & -a_{n-2} \\
& & \ddots & \\
& & & 1 \\
& & & 1 \\
& & -a_{1}
\end{array}\right)
\]

的不变因子组和 Smith 标准形。
\end{document}
